\begin{hyperbolicbox}{Homothetic lattices}
    \small
    We want to regard two lattices to be `similar' if one can be converted to the other via some basic geometric transformations which preserves most of the properties of the lattice structure. Under this idea, we give the following definition.

    \begin{defi}
        Two lattices $\Lambda$ and $\Lambda'$ are said to be \textit{homothetic} if there is $c\in\C\setminus\set{0}$ such that $\Lambda'=c\Lambda$. \cite{milneMF}
    \end{defi}
    
    It's easy to see that, if $\Lambda'=re^{i\theta}\Lambda$ are homothetic, then $\Lambda$ can be transformed to $\Lambda'$ via a dilation by $r$ and a rotation by $\theta$.

    For the structure-preserving property, one has the following theorem.

    \begin{thm}
        Let $\Lambda$ and $\Lambda'$ be two lattices in $\C$. Then $\C/\Lambda$ is isomorphic to $\C/\Lambda'$ if and only if $\Lambda$ and $\Lambda'$ are homothetic. \cite{silvermanTA}
    \end{thm}
\end{hyperbolicbox}